\chapter{Correlation, Association, and the Yule--Simpson Paradox}
\label{chapter::correlationassociation}

\section{Some commonly-used measures of association}

\subsection{Correlation and regression}
\label{section:correlatioandregressionintro}

两个随机变量 $Z$ 与 $Y$ 之间的 \term{Pearson 相关系数}{Pearson correlation coefficient} 定义为
\[
\rho_{ZY}
= \frac{\cov(Z,Y)}{\sqrt{\var(Z)\var(Y)}} ,
\]
它度量 $Z$ 与 $Y$ 的\term{线性依赖}{linear dependence}。

把 $Y$ 对 $Z$ 作\term{线性回归}{linear regression}，得到模型
\begin{equation}
Y = \alpha + \beta Z + \varepsilon ,
\label{eq::short-regression}
\end{equation}
其中 $E(\varepsilon)=0$ 且 $E(\varepsilon Z)=0$。可以证明，\term{回归系数}{regression coefficient} $\beta$ 满足
\[
\beta
= \frac{\cov(Z,Y)}{\var(Z)}
= \rho_{ZY}\sqrt{\frac{\var(Y)}{\var(Z)}} .
\]
因此 $\beta$ 与 $\rho_{ZY}$ 总是同号。

我们也可以定义 $Y$ 对 $Z$ 与 $X$ 的\term{多元回归}{multiple regression}：
\begin{equation}
Y = \alpha + \beta Z + \gamma X + \varepsilon ,
\label{eq::long-regression}
\end{equation}
其中 $E(\varepsilon)=0$，$E(\varepsilon Z)=0$ 且 $E(\varepsilon X)=0$。通常我们会把 $\beta$ 解释为 $Z$ 对 $Y$ 的“效应”：在保持 $X$ 不变、以 $X$ 为条件，或者控制 $X$ 后，$Z$ 与 $Y$ 的关系。线性回归的基础知识见 \cref{appendix::basic-linear-regression}。

更有意思的是，\cref{eq::short-regression,eq::long-regression} 中的两个 $\beta$ 可以不同，甚至可以符号相反。下面的 R 代码重新分析了 \citet{hainmueller2012entropy} 使用的 LaLonde 观测数据。我们关心的主要问题是职业培训项目对收入的“因果效应”。控制所有协变量时，\ri{treat} 的回归系数为 \ri{1067.5461}；不控制任何协变量时，\ri{treat} 的回归系数为 \ri{-8506.4954}。

\begin{lstlisting}
> dat <- read.table("cps1re74.csv", header = TRUE)
> dat$u74 <- as.numeric(dat$re74==0)
> dat$u75 <- as.numeric(dat$re75==0)
>
> ## linear regression on the outcome
> ## . means regression on all other variable in dat
> lmoutcome = lm(re78 ~ ., data = dat)
> round(summary(lmoutcome)$coef[2, ], 3)
  Estimate Std. Error    t value   Pr(>|t|)
  1067.546    554.060      1.927      0.054
>
> lmoutcome = lm(re78 ~ treat, data = dat)
> round(summary(lmoutcome)$coef[2, ], 3)
  Estimate Std. Error    t value   Pr(>|t|)
 -8506.495    712.766    -11.934      0.000
\end{lstlisting}

\subsection{Contingency tables}
\label{sec::twobytwotables-withresume}

两个\term{二元变量}{binary variables} $Z$ 与 $Y$ 的联合分布可以用一个二乘二\term{列联表}{contingency table} 表示。令 $p_{zy}=\pr(Z=z,Y=y)$，则联合分布可以写成
\[
\begin{array}{lcc}
\hline
        & Y=1 & Y=0 \\\hline
Z=1 & p_{11} & p_{10}\\
Z=0 & p_{01} & p_{00}\\
\hline
\end{array}
\]

如果把 $Z$ 看作\term{处理}{treatment}或\term{暴露}{exposure}，把 $Y$ 看作\term{结果}{outcome}，则可以定义\term{风险差}{risk difference, RD}：
\begin{align*}
\RD
&= \pr(Y=1\mid Z=1)-\pr(Y=1\mid Z=0)\\
&= \frac{p_{11}}{p_{11}+p_{10}}
-\frac{p_{01}}{p_{01}+p_{00}},
\end{align*}
\term{风险比}{risk ratio, RR}：
\begin{align*}
\RR
&= \frac{\pr(Y=1\mid Z=1)}{\pr(Y=1\mid Z=0)}\\
&=
\frac{p_{11}}{p_{11}+p_{10}}
\Big/
\frac{p_{01}}{p_{01}+p_{00}},
\end{align*}
以及\term{优势比}{odds ratio, OR}\footnote{在概率论中，一个事件的优势（odds）定义为该事件发生的概率除以该事件不发生的概率。}：
\begin{align*}
\OR
&=
\frac{\pr(Y=1\mid Z=1)/\pr(Y=0\mid Z=1)}
{\pr(Y=1\mid Z=0)/\pr(Y=0\mid Z=0)}\\
&=
\frac{\frac{p_{11}}{p_{11}+p_{10}}/\frac{p_{10}}{p_{11}+p_{10}}}
{\frac{p_{01}}{p_{01}+p_{00}}/\frac{p_{00}}{p_{01}+p_{00}}}\\
&= \frac{p_{11}p_{00}}{p_{10}p_{01}} .
\end{align*}

“风险差”“风险比”和“优势比”这些术语来自\term{流行病学}{epidemiology}。因为流行病学中的结果常常是疾病，所以用“风险”来表示患病概率是自然的。

关于这些度量，有如下简单事实。
\begin{proposition}\label{prop::2x2-independence}
(1) 以下陈述等价\footnote{本书用记号 $\ind$ 表示随机变量的独立或条件独立。这个记号来自 \citet{dawid1979conditional}。}：$Z\ind Y$，$\RD=0$，$\RR=1$，以及 $\OR=1$。

(2) 如果所有 $p_{zy}$ 都为正，则 $\RD>0$ 等价于 $\RR>1$，也等价于 $\OR>1$。

(3) 如果 $\pr(Y=1\mid Z=1)$ 和 $\pr(Y=1\mid Z=0)$ 都很小，则 $\OR\approx \RR$。
\end{proposition}

\cref{prop::2x2-independence} 中 (1) 与 (2) 的证明留作习题 \cref{hw::2x2-equivalence}。(3) 是一个非形式化陈述。近似成立的原因是：当罕见疾病满足 $p\approx 0$ 时，优势 $p/(1-p)$ 接近概率 $p$；由 Taylor 展开，
\[
\frac{p}{1-p}=p+p^2+\cdots\approx p .
\]
在流行病学中，如果结果表示某种罕见疾病的发生，那么假设 $\pr(Y=1\mid X=1)$ 与 $\pr(Y=1\mid X=0)$ 都很小通常是合理的。

如果把概率替换为给定另一个变量 $X$ 后的条件概率，例如 $\pr(Y=1\mid Z=1,X=x)$ 和 $\pr(Y=1\mid Z=0,X=x)$，也可以定义 $\RD$、$\RR$ 和 $\OR$ 的条件版本。

给定计数 $n_{zy}=\#\{i:Z_i=z,Y_i=y\}$，观测数据可以总结为如下二乘二列联表：
\[
\begin{array}{lcc}
\hline
        & Y=1 & Y=0 \\\hline
Z=1 & n_{11} & n_{10}\\
Z=0 & n_{01} & n_{00}\\
\hline
\end{array}
\]
用样本比例 $\hat p_{zy}=n_{zy}/n$ 替换真实概率即可估计 $\RD$、$\RR$ 和 $\OR$，其中 $n$ 是总样本量；二乘二表的渐近推断见 \cref{section::asymptotic-inference-2x2}。在 R 中，\ri{fisher.test} 基于观测数据的二乘二表对 $Z\ind Y$ 作精确检验，\ri{chisq.test} 则作渐近检验。

\begin{example}\label{eg::bertrand2004emily-chapter1}
\citet{bertrand2004emily} 做了一个关于简历的随机实验，研究感知种族对面试回电的影响。他们把听起来像黑人或白人的名字随机分配给虚构简历，再投递给 Boston 和 Chicago 报纸上的招聘广告。下面的二乘二表总结了感知种族与回电结果：

\begin{lstlisting}
> resume = read.csv("resume.csv")
> Alltable = table(resume$race, resume$call)
> Alltable

           0    1
  black 2278  157
  white 2200  235
\end{lstlisting}

两行的总数相同，所以很明显，白人化名字获得了更多回电。下面的 Fisher 精确检验表明，这个差异在统计上显著。
\begin{lstlisting}
> fisher.test(Alltable)

	Fisher's Exact Test for Count Data

data:  Alltable
p-value = 4.759e-05
alternative hypothesis: true odds ratio is not equal to 1
95 percent confidence interval:
 1.249828 1.925573
sample estimates:
odds ratio
  1.549732
\end{lstlisting}
\end{example}

\section{An example of the Yule--Simpson Paradox}

\subsection{Data}

经典的肾结石例子来自 \citet{charig1986comparison}。其中 $Z$ 是处理变量，$Z=1$ 表示开放手术，$Z=0$ 表示小切口手术；$Y$ 是结果变量，$Y=1$ 表示成功，$Y=0$ 表示失败。处理与结果数据可以总结为如下二乘二表：
\[
\begin{array}{lcc}
\hline
        & Y=1 & Y=0 \\\hline
Z=1 & 273 & 77\\
Z=0 & 289 & 61\\
\hline
\end{array}
\]

估计得到的风险差为
\[
\widehat{\RD}
= \frac{273}{273+77}-\frac{289}{289+61}
=78\%-83\%=-5\%<0 .
\]
看起来处理 0 更好；也就是说，相比开放手术，小切口手术有更高的成功率。

不过，这些数据并不是来自\term{随机对照试验}{randomized controlled trial, RCT}\footnote{在 RCT 中，病人被随机分配到不同处理组。本书 \cref{part::rcts} 将集中讨论 RCT。}。接受处理 1 的病人可能和接受处理 0 的病人很不一样。这个研究中有一个\term{潜伏变量}{lurking variable}，即病情严重程度：有些病人的结石较小，有些病人的结石较大。我们可以按结石大小拆分数据。

对结石较小的病人，处理与结果数据为
\[
\begin{array}{lcc}
\hline
        & Y=1 & Y=0 \\\hline
Z=1 & 81 & 6\\
Z=0 & 234 & 36\\
\hline
\end{array}
\]
对结石较大的病人，处理与结果数据为
\[
\begin{array}{lcc}
\hline
        & Y=1 & Y=0 \\\hline
Z=1 & 192 & 71\\
Z=0 & 55 & 25\\
\hline
\end{array}
\]

后两个表相加必须得到第一个表：
\[
81+192=273,\quad
6+71=77,\quad
234+55=289,\quad
36+25=61 .
\]
对结石较小的病人，估计风险差为
\[
\widehat{\RD}_{\text{smaller}}
= \frac{81}{81+6}-\frac{234}{234+36}
=93\%-87\%=6\%>0,
\]
这表明处理 1 更好。对结石较大的病人，估计风险差为
\[
\widehat{\RD}_{\text{larger}}
= \frac{192}{192+71}-\frac{55}{55+25}
=73\%-69\%=4\%>0,
\]
同样表明处理 1 更好。

上面的数据分析得到
\[
\widehat{\RD}<0,\quad
\widehat{\RD}_{\text{smaller}}>0,\quad
\widehat{\RD}_{\text{larger}}>0 .
\]
非正式地说，处理 1 对结石较小和较大的病人都更好，但对总体而言却更差。如果目标是推断处理效应，这个解释就相当令人困惑。在统计学中，这叫作 \term{Yule--Simpson 悖论}{Yule--Simpson Paradox} 或 \term{Simpson 悖论}{Simpson's Paradox}：边际关联的符号与所有层内的条件关联符号相反。

\subsection{Explanation}

令 $X$ 是一个二元指示变量，$X=1$ 表示结石较小，$X=0$ 表示结石较大。先看 $X$ 与 $Z$ 的关系，比较小结石病人与大结石病人接受处理 1 的概率：
\begin{align*}
&\widehat{\pr}(Z=1\mid X=1)-\widehat{\pr}(Z=1\mid X=0)\\
&=
\frac{81+6}{81+6+234+36}
-\frac{192+71}{192+71+55+25}\\
&= 24\%-77\%\\
&= -53\%<0 .
\end{align*}
因此，大结石病人更经常接受处理 1。从统计上说，$X$ 与 $Z$ 存在负关联。

再看 $X$ 与 $Y$ 的关系，比较小结石病人与大结石病人的成功概率。在处理 1 下，
\begin{align*}
&\widehat{\pr}(Y=1\mid Z=1,X=1)
-\widehat{\pr}(Y=1\mid Z=1,X=0)\\
&= \frac{81}{81+6}-\frac{192}{192+71}\\
&= 93\%-73\%\\
&=20\%>0;
\end{align*}
在处理 0 下，
\begin{align*}
&\widehat{\pr}(Y=1\mid Z=0,X=1)
-\widehat{\pr}(Y=1\mid Z=0,X=0)\\
&= \frac{234}{234+36}-\frac{55}{55+25}\\
&= 87\%-69\%\\
&=18\%>0 .
\end{align*}
因此，在两个处理水平下，小结石病人的成功概率都更高。从统计上说，在两个处理水平内，$X$ 与 $Y$ 都存在正关联。

\begin{figure}[H]
\[
\xymatrix{
 & X\ar_{-}[ld]\ar^{+}[rd] \\
Z \ar_{+}[rr] &&Y
}
\]
\caption{肾结石例子的示意图。符号表示两个变量之间的关联方向，其中会以所有指向下游变量的其他变量为条件。}
\label{fig::kidney-stone}
\end{figure}

我们可以用 \cref{fig::kidney-stone} 总结这些定性关联。用技术语言说，处理到结果有一条正路径 $(Z\rightarrow Y)$，同时还有一条更强的负路径 $(Z\leftarrow X\rightarrow Y)$，所以处理与结果之间的总体关联为负。用普通语言说，当效果较差的处理 0 更常用于病情较轻的病例时，它可能看起来像是更有效的处理。

一般来说，因为 $X$ 与 $Z$ 之间有关联，且 $X$ 与 $Y$ 之间也有关联，$Z$ 与 $Y$ 的关联可能在边际上和给定 $X$ 后的条件层内呈现不同的定性方向。当 Yule--Simpson 悖论发生时，我们说 $X$ 是 $Z$ 与 $Y$ 关系中的\term{混杂变量}{confounding variable}或\term{混杂因子}{confounder}，也可以说 $Z$ 与 $Y$ 的关系被 $X$ \term{混杂}{confounded}。更深入的讨论见 \cref{chapter::observational-studies,chapter::evalue}。

\subsection{Geometry of the Yule--Simpson Paradox}

假设基于汇总数据的二乘二表有如下计数：
\[
\begin{array}{lcc}
\hline
\text{总体} & Y=1 & Y=0\\\hline
Z=1 & n_{11} & n_{10}\\
Z=0 & n_{01} & n_{00}\\
\hline
\end{array}
\]

基于亚组的两个二乘二表分别为
\[
\begin{array}{lcc}
\hline
\text{子总体 } X=1 & Y=1 & Y=0\\\hline
Z=1 & n_{11|1} & n_{10|1}\\
Z=0 & n_{01|1} & n_{00|1}\\
\hline
\end{array}
\]
以及
\[
\begin{array}{lcc}
\hline
\text{子总体 } X=0 & Y=1 & Y=0\\\hline
Z=1 & n_{11|0} & n_{10|0}\\
Z=0 & n_{01|0} & n_{00|0}\\
\hline
\end{array}
\]

\cref{fig::geomtry-simpson} 展示了 Yule--Simpson 悖论的几何图像。纵轴表示 $Y=1$ 的成功次数，横轴表示 $Y=0$ 的失败次数。两个平行四边形对应于在两个处理水平下汇总成功和失败的计数。$OA_1$ 的斜率大于 $OB_1$ 的斜率，$OA_0$ 的斜率也大于 $OB_0$ 的斜率，所以在 $X$ 的两个水平内，处理看起来都对结果有益。然而，$OA$ 的斜率小于 $OB$ 的斜率，所以对总体而言，处理看起来反而有害。Yule--Simpson 悖论由此产生。

\begin{figure}[H]
\includegraphics[width=\textwidth]{figures/simpson_geometry2.pdf}
\caption{Yule--Simpson 悖论的几何图像}
\label{fig::geomtry-simpson}
\end{figure}

\section{The Berkeley graduate school admission data}
\label{sec::berkeleydata}

\citet{bickel1975sex} 研究了 Berkeley 研究生院男女学生的录取率。R 的 \ri{datasets} 包含原始数据 \ri{UCBAdmissions}。六个最大院系的原始数据如下：
\begin{lstlisting}
> library(datasets)
> UCBAdmissions = aperm(UCBAdmissions, c(2, 1, 3))
> UCBAdmissions
, , Dept = A

        Admit
Gender   Admitted Rejected
  Male        512      313
  Female       89       19

, , Dept = B

        Admit
Gender   Admitted Rejected
  Male        353      207
  Female       17        8

, , Dept = C

        Admit
Gender   Admitted Rejected
  Male        120      205
  Female      202      391

, , Dept = D

        Admit
Gender   Admitted Rejected
  Male        138      279
  Female      131      244

, , Dept = E

        Admit
Gender   Admitted Rejected
  Male         53      138
  Female       94      299

, , Dept = F

        Admit
Gender   Admitted Rejected
  Male         22      351
  Female       24      317
\end{lstlisting}

把数据按院系汇总后，得到一个简单的二乘二表：
\begin{lstlisting}
> UCBAdmissions.sum = apply(UCBAdmissions, c(1, 2), sum)
> UCBAdmissions.sum
        Admit
Gender   Admitted Rejected
  Male       1198     1493
  Female      557     1278
\end{lstlisting}

下面这个函数基于 \ri{chisq.test}，输入一个二乘二表，输出估计的风险差和 $p$ 值：
\begin{lstlisting}
> risk.difference = function(tb2)
+ {
+   p1      = tb2[1, 1]/(tb2[1, 1] + tb2[1, 2])
+   p2      = tb2[2, 1]/(tb2[2, 1] + tb2[2, 2])
+   testp   = chisq.test(tb2)
+
+   return(list(p.diff = p1 - p2,
+               pv = testp$p.value))
+ }
\end{lstlisting}
利用这个函数，可以发现男女学生录取率之间存在很大且显著的差异：
\begin{lstlisting}
> risk.difference(UCBAdmissions.sum)
$p.diff
[1] 0.1416454

$pv
[1] 1.055797e-21
\end{lstlisting}

按院系分层后，男女学生录取率的差异变得更小且不显著。在 A 院系中，差异显著但为负。
\begin{lstlisting}
> P.diff = rep(0, 6)
> PV     = rep(0, 6)
> for(dd in 1:6)
+ {
+ 	 department = risk.difference(UCBAdmissions[, , dd])
+ 	 P.diff[dd] = department$p.diff
+ 	 PV[dd]     = department$pv
+ }
>
> round(P.diff, 2)
[1] -0.20 -0.05  0.03 -0.02  0.04 -0.01
> round(PV, 2)
[1] 0.00 0.77 0.43 0.64 0.37 0.64
\end{lstlisting}

\section{Homework Problems}

\homeworksolutionref{01}
\paragraph{Independence in two-by-two tables}\label{hw::2x2-equivalence}

证明 \cref{prop::2x2-independence} 中的 (1) 与 (2)。

\paragraph{More examples of the Yule--Simpson Paradox}

给出一个数值例子，构造一个会产生 Yule--Simpson 悖论的二乘二乘二表。

找一个现实生活中出现 Yule--Simpson 悖论的例子。

\paragraph{Correlation and partial correlation}\label{hw::partial-correlation}

考虑一个三维正态随机向量：
\[
\begin{pmatrix}
X\\
Y\\
Z
\end{pmatrix}
\sim
\textsc{N}
\left(
\begin{pmatrix}
0\\
0\\
0
\end{pmatrix},
\begin{pmatrix}
1 & \rho_{XY} & \rho_{XZ}\\
\rho_{XY} & 1 & \rho_{YZ}\\
\rho_{XZ} & \rho_{YZ} & 1
\end{pmatrix}
\right).
\]
$Y$ 与 $Z$ 的相关系数是 $\rho_{YZ}$。偏相关系数有许多等价定义。对一个多元正态向量，令 $\rho_{YZ\mid X}$ 表示给定 $X$ 后 $Y$ 与 $Z$ 的偏相关系数，即条件分布 $(Y,Z)\mid X$ 中二者的相关系数。证明
\[
\rho_{YZ\mid X}
=
\frac{\rho_{YZ}-\rho_{YX}\rho_{ZX}}
{\sqrt{1-\rho_{YX}^2}\sqrt{1-\rho_{ZX}^2}} .
\]

给出一个数值例子，使得 $\rho_{YZ}>0$ 且 $\rho_{YZ\mid X}<0$。

注：这是正态随机向量中的 Yule--Simpson 悖论。可以使用 \cref{sec::multivariate-normal} 的结果证明偏相关系数公式。

\paragraph{Specification searches}\label{hw::lalonde-obs-reg}

\cref{section:correlatioandregressionintro} 重新分析了 \citet{hainmueller2012entropy} 使用的数据。我们使用名为 \ri{re78} 的结果变量和名为 \ri{treat} 的处理变量。此外，数据还包含 $10$ 个协变量，因此线性回归中有 $2^{10}=1024$ 个可能的协变量子集。请运行全部 $1024$ 个线性回归，并报告处理变量的回归系数。有多少系数显著为正，有多少显著为负，有多少不显著？也可以报告这些回归中的其他有趣发现。

\paragraph{More on racial discrimination}

\cref{sec::twobytwotables-withresume} 重新分析了 \citet{bertrand2004emily} 收集的数据。请分别对男性和女性做分析。你有什么发现？
